% libre-compendium
%
% Copyright (C) 2026 Richard Qin
%
% This document is dual-licensed under either:
% - GNU Free Documentation License, Version 1.3 or any later version
%   published by the Free Software Foundation; with no Invariant Sections,
%   no Front-Cover Texts, and no Back-Cover Texts (see COPYING.DOC); or
% - Creative Commons Attribution-ShareAlike 4.0 International License
%   (see LICENSE-CC-BY-SA-4.0).

% Exercises migrated from problems02.tex and solutions02.tex.
\begin{exercise}[level=basic]
判断下列映射是否线性，并说明理由：
\begin{enumerate}[label=(\arabic*)]
\item $T:\R^2\to\R^2$，$T(x_1,x_2)=(x_1+2x_2,3x_1)$；
\item $T:\R^2\to\R^2$，$T(x_1,x_2)=(x_1+1,x_2)$；
\item $D:P_3(\R)\to P_2(\R)$，$D(f)=f'$；
\item $T:M_2(\R)\to M_2(\R)$，$T(A)=A+A^{\mathsf T}$。
\end{enumerate}
\begin{solution}
(1) 线性，它由矩阵 $\begin{pmatrix}1&2\\3&0\end{pmatrix}$ 给出。(2) 不线性，因为 $T(0)=(1,0)\ne0$。(3) 线性，微分保持加法和数乘。(4) 线性，因为
\[
T(aA+bB)=a(A+A^{\mathsf T})+b(B+B^{\mathsf T}).
\]
\end{solution}
\end{exercise}

\begin{exercise}[level=basic]
设 $T:\R^3\to\R^2$，
\[
T(x_1,x_2,x_3)=(x_1+x_2,\ x_2+x_3).
\]
求 $\ker T$、$\operatorname{Im}T$、秩与零化度。
\begin{solution}
核由
\[
x_1+x_2=0,\qquad x_2+x_3=0
\]
确定，故
\[
\ker T=\operatorname{span}\{(-1,1,-1)\}.
\]
矩阵 $\begin{pmatrix}1&1&0\\0&1&1\end{pmatrix}$ 行满秩，所以
\[
\operatorname{Im}T=\R^2,\quad
\rank T=2,\quad
\operatorname{nullity}T=1.
\]
\end{solution}
\end{exercise}

\begin{exercise}[level=basic]
设 $T:V\to W$ 线性，$T(x_0)=y$。证明 $T(x)=y$ 的全部解为 $x_0+\ker T$。
\begin{solution}
若 $z\in\ker T$，则 $T(x_0+z)=y$。若 $T(x)=y$，则
\[
T(x-x_0)=0,
\]
故 $x-x_0\in\ker T$。两边包含关系均成立。
\end{solution}
\end{exercise}

\begin{exercise}[level=basic]
设 $(e_1,e_2,e_3)$ 是 $V$ 的一组基，指定
\[
T(e_1)=e_1+e_2,\quad
T(e_2)=e_2+e_3,\quad
T(e_3)=e_1-e_3.
\]
求 $T(2e_1-e_2+3e_3)$ 以及 $T$ 在该基下的矩阵。
\begin{solution}
由线性性，
\begin{align*}
T(2e_1-e_2+3e_3)
&=2(e_1+e_2)-(e_2+e_3)+3(e_1-e_3)\\
&=5e_1+e_2-4e_3.
\end{align*}
矩阵各列为三个基向量的像坐标：
\[
[T]=
\begin{pmatrix}
1&0&1\\
1&1&0\\
0&1&-1
\end{pmatrix}.
\]
\end{solution}
\end{exercise}

\begin{exercise}[level=basic]
设
\[
A=\begin{pmatrix}1&2&0\\0&1&1\end{pmatrix}.
\]
把 $A$ 看作线性映射 $T_A:\R^3\to\R^2$ 的矩阵。利用秩零化度定理求 $\dim\ker T_A$，并求核的一组基。
\begin{solution}
$A$ 的秩为 $2$，故
\[
\dim\ker T_A=3-2=1.
\]
解方程
\[
x_1+2x_2=0,\qquad x_2+x_3=0
\]
得
\[
(x_1,x_2,x_3)=t(-2,1,-1).
\]
\end{solution}
\end{exercise}

\begin{exercise}[level=basic]
在 $\R^3$ 中取
\[
V_1=\operatorname{span}\{(1,0,1),(0,1,1)\},
\quad
V_2=\operatorname{span}\{(1,1,2),(1,-1,0)\}.
\]
求 $V_1+V_2$ 与 $V_1\cap V_2$ 的维数。
\begin{solution}
注意
\[
(1,1,2)=(1,0,1)+(0,1,1),
\]
\[
(1,-1,0)=(1,0,1)-(0,1,1).
\]
所以 $V_2=V_1$。两个生成向量线性无关，因此
\[
\dim(V_1+V_2)=2,\qquad
\dim(V_1\cap V_2)=2.
\]
\end{solution}
\end{exercise}

\begin{exercise}[level=basic]
证明：若 $V=V_1\oplus V_2$，则 $V_1$ 的一组基与 $V_2$ 的一组基合并后构成 $V$ 的一组基。
\begin{solution}
设 $(u_1,\ldots,u_r)$、$(w_1,\ldots,w_s)$ 分别为两子空间的基。直和保证每个 $v\in V$ 可写成 $u+w$，所以合并向量组张成 $V$。若
\[
\sum a_iu_i+\sum b_jw_j=0,
\]
直和的零表示唯一性给出两部分均为零，再由各自基的线性无关性得到所有系数为零。
\end{solution}
\end{exercise}

\begin{exercise}[level=standard]
设 $T:P_3(\R)\to\R^3$，
\[
T(f)=(f(0),f(1),f'(0)).
\]
求核、像、秩与零化度，并判断 $T$ 是否满射。
\begin{solution}
设 $f=a+bt+ct^2+dt^3$，则
\[
T(f)=(a,a+b+c+d,b).
\]
核条件给出 $a=b=0$ 与 $c+d=0$，故
\[
\ker T=\operatorname{span}\{t^2-t^3\}.
\]
零化度为 $1$，秩零化度定理给出秩为 $3$。陪域也是三维，所以像等于 $\R^3$，$T$ 满射。
\end{solution}
\end{exercise}

\begin{exercise}[level=standard]
设 $T:\R^2\to\R^2$ 在标准基下矩阵为
\[
A=\begin{pmatrix}2&1\\0&3\end{pmatrix}.
\]
取新基
\[
\widetilde{\mathcal E}
=((1,1)^{\mathsf T},(1,-1)^{\mathsf T}).
\]
求 $T$ 在新基下的矩阵。
\begin{solution}
过渡矩阵
\[
P=\begin{pmatrix}1&1\\1&-1\end{pmatrix},
\qquad
P^{-1}=\frac12\begin{pmatrix}1&1\\1&-1\end{pmatrix}.
\]
新基下矩阵为
\[
P^{-1}AP
=\begin{pmatrix}
3&-1\\
0&2
\end{pmatrix}.
\]
\end{solution}
\end{exercise}

\begin{exercise}[level=standard]
设 $T:P_2(\R)\to P_2(\R)$，$T(f)=f+f'$。求它在基 $(1,t,t^2)$ 下的矩阵，并判断是否可对角化。
\begin{solution}
\[
T(1)=1,\qquad
T(t)=1+t,\qquad
T(t^2)=2t+t^2.
\]
故
\[
[T]=\begin{pmatrix}
1&1&0\\
0&1&2\\
0&0&1
\end{pmatrix}.
\]
唯一特征值为 $1$，但
\[
\dim\ker([T]-I)=1<3,
\]
所以不可对角化。
\end{solution}
\end{exercise}

\begin{exercise}[level=standard]
设 $A$ 的特征多项式为
\[
(\lambda-1)^3(\lambda+2)^2,
\]
且
\[
\dim\ker(A-I)=2,\qquad
\dim\ker(A+2I)=2.
\]
判断 $A$ 是否可对角化。
\begin{solution}
特征值 $1$ 的几何重数为 $2<3$，没有达到代数重数，因此 $A$ 不可对角化。虽然特征值 $-2$ 的几何重数已达到 $2$，但必须每个特征值都满足条件。
\end{solution}
\end{exercise}

\begin{exercise}[level=standard]
证明：不同特征值对应的特征向量线性无关。
\begin{solution}
对向量个数归纳。设
\[
a_1v_1+\cdots+a_sv_s=0,
\qquad
T(v_i)=\lambda_iv_i,
\]
其中 $\lambda_i$ 两两不同。对等式作用 $T-\lambda_sI$，得到
\[
\sum_{i=1}^{s-1}a_i(\lambda_i-\lambda_s)v_i=0.
\]
由归纳假设，$a_i(\lambda_i-\lambda_s)=0$，故 $a_1=\cdots=a_{s-1}=0$，原式再给出 $a_s=0$。
\end{solution}
\end{exercise}

\begin{exercise}[level=standard]
判断以下三个子空间之和是否为直和：
\[
V_1=\operatorname{span}\{(1,0,0)\},\quad
V_2=\operatorname{span}\{(0,1,0)\},\quad
V_3=\operatorname{span}\{(1,1,1)\}.
\]
若是，求其和空间。
\begin{solution}
三个生成向量组成矩阵
\[
\begin{pmatrix}
1&0&1\\
0&1&1\\
0&0&1
\end{pmatrix},
\]
行列式为 $1$，故线性无关。三个一维子空间之和是直和，并且等于 $\R^3$。
\end{solution}
\end{exercise}

\begin{exercise}[level=standard]
设 $P:V\to V$ 满足 $P^2=P$。证明
\[
\ker P=\operatorname{Im}(I-P),
\qquad
\operatorname{Im}P=\ker(I-P).
\]
\begin{solution}
若 $x\in\ker P$，则 $(I-P)x=x$，故 $x\in\operatorname{Im}(I-P)$；反之，
\[
P(I-P)=P-P^2=0,
\]
故 $\operatorname{Im}(I-P)\subseteq\ker P$。第二式由
\[
(I-P)P=P-P^2=0
\]
以及 $x\in\ker(I-P)\Rightarrow x=Px$ 同理得到。
\end{solution}
\end{exercise}

\begin{exercise}[level=standard]
在 $M_n(\R)$ 上定义
\[
P(A)=\frac{A+A^{\mathsf T}}2.
\]
证明 $P$ 是投影，求其像与核。
\begin{solution}
\[
P^2(A)
=\frac12\left(P(A)+P(A)^{\mathsf T}\right)
=P(A),
\]
因为 $P(A)$ 已经对称。所以 $P$ 是投影。其像是全体对称矩阵；核由
\[
A+A^{\mathsf T}=0
\]
确定，是全体反对称矩阵。
\end{solution}
\end{exercise}

\begin{exercise}[level=advanced]
设 $T:V\to W$、$S:W\to U$ 为有限维线性映射。证明
\[
\rank(S\circ T)\le
\min\{\rank S,\rank T\}.
\]
\begin{solution}
\[
\operatorname{Im}(S\circ T)=S(\operatorname{Im}T)
\subseteq\operatorname{Im}S,
\]
故复合秩不超过 $\rank S$。同时 $S$ 在线性空间 $\operatorname{Im}T$ 上的像维数不超过定义域维数，所以
\[
\rank(S\circ T)\le\dim\operatorname{Im}T=\rank T.
\]
\end{solution}
\end{exercise}

\begin{exercise}[level=advanced]
设 $T:V\to W$ 线性。证明：
\[
\dim T(M)=\dim M-\dim(M\cap\ker T)
\]
对任意有限维子空间 $M\subseteq V$ 成立。
\begin{solution}
考虑限制映射
\[
T|_M:M\to W.
\]
其核为
\[
\ker(T|_M)=M\cap\ker T,
\]
像为 $T(M)$。对限制映射使用秩零化度定理：
\[
\dim M=\dim T(M)+\dim(M\cap\ker T).
\]
\end{solution}
\end{exercise}

\begin{exercise}[level=advanced]
设 $T:V\to V$ 满足 $T^2=T$。证明其最小多项式整除 $x(x-1)$，并由此证明 $T$ 可对角化。
\begin{solution}
由 $T^2=T$，
\[
T(T-I)=0,
\]
所以多项式 $x(x-1)$ 消去 $T$，最小多项式整除它。多项式 $x(x-1)$ 在任意数域上分裂且无重根，因此最小多项式也分裂且无重根，故 $T$ 可对角化。也可直接使用
\[
V=\operatorname{Im}T\oplus\ker T
\]
并在两部分上分别得到特征值 $1,0$。
\end{solution}
\end{exercise}

\begin{exercise}[level=advanced]
设 $T:V\to V$，且
\[
V=V_1\oplus V_2,
\]
其中 $V_1,V_2$ 都是 $T$-不变子空间。证明
\[
\chi_T(\lambda)
=\chi_{T|_{V_1}}(\lambda)
\chi_{T|_{V_2}}(\lambda).
\]
\begin{solution}
分别取 $V_1,V_2$ 的基并合并。矩阵为
\[
[T]=\begin{pmatrix}A_1&0\\0&A_2\end{pmatrix}.
\]
因此
\[
\chi_T(\lambda)
=\det(\lambda I-A_1)\det(\lambda I-A_2)
=\chi_{T|_{V_1}}(\lambda)\chi_{T|_{V_2}}(\lambda).
\]
\end{solution}
\end{exercise}

\begin{exercise}[level=advanced]
设 $V$ 有限维，$V=U\oplus W$。定义
\[
P(u+w)=u.
\]
证明 $P$ 是线性映射，并证明满足
\[
\operatorname{Im}P=U,\qquad
\ker P=W
\]
的投影算子唯一。
\begin{solution}
若 $x_i=u_i+w_i$，则
\[
P(ax_1+bx_2)=a u_1+b u_2=aP(x_1)+bP(x_2),
\]
故 $P$ 线性。由定义直接得到像为 $U$、核为 $W$。

若 $Q$ 是另一个具有相同像、核且为投影的算子，对 $x=u+w$，由于 $u\in\operatorname{Im}Q$，有 $Q(u)=u$；由于 $w\in\ker Q$，有 $Q(w)=0$。故 $Q(x)=u=P(x)$。
\end{solution}
\end{exercise}

\begin{exercise}[level=advanced]
设 $V_1,V_2,V_3$ 是有限维空间 $V$ 的子空间，且两两交为零。给出它们的和不为直和的例子，并给出保证
\[
V_1+V_2+V_3
\]
为直和的正确充要条件。
\begin{solution}
在 $\R^2$ 中取
\[
V_1=\operatorname{span}\{e_1\},\quad
V_2=\operatorname{span}\{e_2\},\quad
V_3=\operatorname{span}\{e_1+e_2\}.
\]
它们两两交为零，但存在非平凡关系
\[
e_1+e_2-(e_1+e_2)=0.
\]
正确充要条件是：若 $x_i\in V_i$ 且 $x_1+x_2+x_3=0$，则所有 $x_i=0$；等价地，
\[
V_i\cap\sum_{j\ne i}V_j=\{0\}
\]
对每个 $i$ 成立。
\end{solution}
\end{exercise}

\begin{exercise}[level=advanced]
设 $T:\R^2\to\R^3$，
\[
T(x_1,x_2)=(x_1,x_2,x_1+x_2).
\]
构造一个左逆 $L:\R^3\to\R^2$。该映射是否存在右逆？
\begin{solution}
$T$ 单射。取
\[
L(y_1,y_2,y_3)=(y_1,y_2),
\]
则
\[
(L\circ T)(x_1,x_2)=(x_1,x_2),
\]
所以 $L$ 是左逆。$T$ 不满射，因为其像为平面
\[
y_3=y_1+y_2,
\]
故不存在右逆。
\end{solution}
\end{exercise}

\begin{exercise}[level=advanced]
设 $A\in\F^{m\times n}$、$B\in\F^{n\times p}$。利用 Sylvester 秩不等式证明：若 $AB=0$，则
\[
\rank A+\rank B\le n.
\]
\begin{solution}
Sylvester 秩不等式给出
\[
\rank(AB)\ge\rank A+\rank B-n.
\]
由 $AB=0$，左边为零，因此
\[
\rank A+\rank B\le n.
\]
\end{solution}
\end{exercise}

\begin{exercise}[level=advanced]
设 $\dim V=3$、$\dim W=2$。求 $\dim\operatorname{Hom}(V,W)$，并说明如何由两组基构造其一组基。
\begin{solution}
\[
\dim\operatorname{Hom}(V,W)
=(\dim V)(\dim W)=6.
\]
取 $V$ 的基 $(e_1,e_2,e_3)$ 和 $W$ 的基 $(f_1,f_2)$。对每个 $(i,j)$ 定义 $E_{ij}:V\to W$，使
\[
E_{ij}(e_j)=f_i,
\]
并在其余基向量上取零。六个映射 $E_{ij}$ 对应六个矩阵单位，构成一组基。
\end{solution}
\end{exercise}

\begin{exercise}[level=advanced]
设
\[
A=\begin{pmatrix}1&1\\0&2\end{pmatrix}.
\]
求其最小多项式，并利用最小多项式判据说明 $A$ 可对角化。
\begin{solution}
特征多项式为
\[
\chi_A(x)=(x-1)(x-2).
\]
$A$ 既不等于 $I$ 也不等于 $2I$，故最小多项式不可能是一次因子；又由 Cayley--Hamilton 定理，它整除特征多项式，所以
\[
m_A(x)=(x-1)(x-2).
\]
该多项式在 $\R$ 上分裂且无重根，因此 $A$ 可对角化。
\end{solution}
\end{exercise}

\begin{exercise}[level=advanced]
对上题矩阵，利用谱投影公式求 $P_1,P_2$，验证
\[
P_1+P_2=I,\qquad
A=P_1+2P_2.
\]
\begin{solution}
对特征值 $1,2$，
\[
P_1=\frac{A-2I}{1-2}=2I-A
=\begin{pmatrix}1&-1\\0&0\end{pmatrix},
\]
\[
P_2=\frac{A-I}{2-1}=A-I
=\begin{pmatrix}0&1\\0&1\end{pmatrix}.
\]
直接相加得 $P_1+P_2=I$，并且
\[
P_1+2P_2
=\begin{pmatrix}1&1\\0&2\end{pmatrix}=A.
\]
还可验证 $P_i^2=P_i$ 与 $P_1P_2=0$。
\end{solution}
\end{exercise}

\begin{exercise}[level=advanced]
设
\[
J=\begin{pmatrix}1&1\\0&1\end{pmatrix},
\qquad
M=\operatorname{span}\{(1,0)^{\mathsf T}\}.
\]
证明 $M$ 是 $J$-不变子空间，但不存在 $J$-不变子空间 $N$ 使 $\R^2=M\oplus N$。
\begin{solution}
\[
J(1,0)^{\mathsf T}=(1,0)^{\mathsf T},
\]
故 $M$ 不变。若存在一维不变补空间 $N$，其任一非零生成向量必须是 $J$ 的特征向量。唯一特征值为 $1$，而
\[
\ker(J-I)=\operatorname{span}\{(1,0)^{\mathsf T}\}=M.
\]
所以所有特征向量都属于 $M$，不能张成与 $M$ 互补的直线，矛盾。
\end{solution}
\end{exercise}

